%\vspace{-0.5em}
\section{Related Works}
%\vspace{-0.3em}
\noindent\textbf{Intent Analysis and Understanding: }
%
Several papers dealt with the concept of
intents or goals for recommendations in process mining, web mining, education~\cite{jiang2019time, jiang2019goal} and analyze patterns in user behavior from application log/clicks data~\cite{aggarwal2020goal, dev2017identifying} or characterised SQL query logs to understand analysis behaviors~\cite{jain2016sqlshare}. 
%They do not consider the results for each of the those queries that might have influenced the patterns.
Distinct from all these, we identify goal as a combination of queries and the corresponding results in a sequential decision making workflow.

\noindent\textbf{Data Exploration:}
%
\cite{brachmann2019data} simplified
the manual creation EDA notebooks, cite{chirigati2016data} focused on related dataset recommendation. \cite{zhao2017controlling} attempts to flag false discoveries in EDA automation but does not consider approximations or intent of analysts.
%
Recent work also focused on automatically generating  data insights~\cite{ma2021metainsight, ding2019quickinsights} or creating or augmenting EDA workflows~\cite{atena2020sigmod, kery2018story, rule2018exploration}. 
%As an orthogonal direction, we focus on faster exploration with intent-aware use of approximations.
Several papers used samples and approximate query processing to speed up data explorations and analytics~\cite{sheoran2022aaai, sheoran2022electra, babcock2003dynamic, chaudhuri2017approximate, ma2021learned, bater2020saqe} but we are the first to look at the sequence and context of such exploration for visual analytics. 

%\noindent\textbf{Sampling and Approximations:} 
%
%\cite{park2016visualization} proposed a sampling strategy that guarantees high quality visualizations with a small subset of the entire dataset. But they do not consider interactive and sequential decision making workflow. 
%Several works~\cite{moritz2017trust, zhao2018random, agarwal2013blinkdb, park2018verdictdb} developed techniques for sampling to reduce error for ad-hoc queries, and make them compute efficient~\cite{ding2016sample, peng2018aqp}. We used several of these sampling techniques in our action space.

\noindent\textbf{Interactive Applications by Reinforcement Learning}
%The recent advances of reinforcement learning and deep learning enables many novel interactive applications.
\citet{liu2018dialogue,takanobu2020multi} propose interactive dialog systems based on reinforcement learning.
\citet{guo2018dialog} propose a dialog-based interactive image retrieval system involving natural language feedback on visual attributes of the items.
%Instead of searching for simple images,
\citet{tan2019drill} propose reinforcement learning with a drill-down method to interactively retrieve complex scene images.
Reinforcement learning from reformulations is proposed for question answering, by modeling the answering process as multiple agents walking in parallel on the knowledge graph \cite{kaiser2021reinforcement}.  \citet{hua2020few} propose
a meta reinforcement learning method in complex question answering, which quickly adapts to unseen questions. 
Reinforcement learning has also been applied to elicit user preferences in conversational recommender systems \cite{christakopoulou2016towards, zhang2020conversational,lei2020interactive}. Similar to these interactive scenarios \cite{guo2018dialog,yu2019visual,takanobu2020multi}, a major challenge in the evaluation is that we need to have access to user reactions to any possible results returned by the system, which are exhaustive to collect. Therefore, similarly, we adopt a setting where the user reactions are generated by simulators as surrogates for real users as detailed in Section \ref{sec:evaluation}.

%\noindent\textbf{Temporal Question Answering}
%\sm{Missing discussion of the related work in “temporal question answering”}
%\tong{not sure whether we need a separate section}
%\tong{Please summarize some relevant papers.}
%\sm{Tong: Please write about use of RL and simulators in a wide range of domains like image captioning, NLP etc.}

